\begin{table}[H]
\centering
\caption{Lee (2009) Trimming Bounds (Ages 10--13)}
\label{tab:lee_bounds}
\begin{threeparttable}
\begin{tabular}{lccccc}
\toprule
Subgroup & Point Est. & Lower Bound & Upper Bound & Trim \% & Selection $p$ \\
\midrule
Full sample (ages 10-13) & -0.0056 & -0.0056 & -0.0056 & 0.0 & 0.003 \\
Full sample: enrollment & 0.0161 & 0.0161 & 0.0161 & 0.0 & 0.651 \\
Girls: child labor & -0.0194 & -0.0194 & -0.0194 & 0.0 & 0.000 \\
Girls: enrollment & 0.0325 & 0.0325 & 0.0325 & 0.0 & 0.779 \\
Middle+Richer: enrollment & 0.0506 & 0.0506 & 0.0506 & 0.0 & 0.083 \\
\bottomrule
\end{tabular}
\begin{tablenotes}
\footnotesize
\item Notes: Lee (2009) trimming bounds account for potential selective attrition. Point estimates from standard DiD with controls. Lower and upper bounds trim the ``excess observations from the treatment-period cell with higher response rate, from the top and bottom of the outcome distribution respectively. Trim \% indicates the fraction of observations trimmed. Selection $p$ tests for differential non-response ($H_0$: treatment does not affect observation probability). All specifications include age dummies, household size, and wealth quintile (or sex where applicable). Standard errors clustered at district level.
\end{tablenotes}
\end{threeparttable}
\end{table}